\section{The decentralized economy}\label{sec:mkt:setup}

This part takes the tastes and technologies of Section~\ref{sec:sp:setup} exactly as they stand and
asks a different question of them: what happens when the allocation is made by competitive firms and
a representative household rather than by a planner, and what the government has to do to make the
two coincide. Nothing in Section~\ref{sec:sp:setup} is modified. What changes is who chooses what,
and --- crucially --- what each chooser is able to see.

\subsection{What the market can and cannot see}\label{sec:mkt:setup:externalities}

It is worth being explicit at the outset about where the decentralization is going to bite, because
the answer is not the one the environmental literature usually starts from. Three distinct failures
are at work.

The first is the familiar one. Pollution is a stock externality: $P$ enters $u$ and $F$, and no
private agent perceives that its own waste adds to it. In the planner's problem $P$ carries a costate
$p^P$; in the market economy no private agent carries a costate on $P$ at all, because $P$ is not a
private state variable. This is the gap that a Pigouvian tax closes, and it is the only one that
appears in the simpler model of \textcite{sorensen2018}.

The second is specific to the materials accounting. The ledger~\eqref{eq:sp:ledger} says that every
tonne entering production must eventually leave the economy as waste. That is a constraint on the
economy, not on any firm: the final-goods producer who buys a tonne of $R$ does not thereby acquire an
obligation to dispose of a tonne later, because the tonne leaves through other agents' hands ---
through the consumer who uses up the good, through the depreciation of somebody else's capital. The
shadow price $p^W$ attached to the ledger therefore has no market counterpart unless one is created.

The third is the subtlest and has no analogue in the earlier literature at all. The intensity
condition~\eqref{eq:sp:intensity} determines $\Omega$, and $\Omega$ is an economy-wide object: each
agent takes the waste intensity of its own activity as a technological datum, while collectively
their choices determine it. What the composition of final-good uses does to the allocation --- the
whole family of wedges $1\pm\zeta\Omega^j$ derived in Section~\ref{sec:sp:opt:foc} --- is therefore
invisible to every individual agent even though it is a real social cost. The multiplier $\zeta$
prices it, and a second instrument is needed to make anyone pay it.

The three failures are logically independent, and the instrument set below has three corresponding
blocks. That the second and third can be handled by prices at all --- rather than by quantity
regulation of $\Omega$, which no government could observe --- is the substantive content of
Section~\ref{sec:mkt:implementation}.

\subsection{Agents, markets and prices}\label{sec:mkt:setup:agents}

The final good is the numéraire. Five kinds of agent trade in four markets.

\smalltitle{Agents.} A representative \emph{household} owns the capital stock and all firms, consumes,
and saves. A competitive \emph{final-goods sector} rents capital and buys raw material. A competitive
\emph{mining sector} extracts virgin material from a reserve it owns and explores for new reserves. A
competitive \emph{recycling sector} rents capital and buys treated waste. A \emph{waste-management
sector} operates under a public collection obligation: it must collect every tonne of waste generated,
it chooses how much of it to treat, and it sells the treated fraction to recyclers. The right to
operate it is awarded by competitive tender, which drives its profit to zero.

\smalltitle{Markets and prices.} Throughout, bare $P$ remains the pollution stock of
Section~\ref{sec:sp:setup}; every market price carries a superscript.
\begin{align}
P^R \;\text{[raw material]},
&&
P^T \;\text{[treated waste]},
&&
P^K \;\text{[capital rental]},
&&
1 \;\text{[final good]}.
\label{eq:mkt:prices}
\end{align}
Virgin and secondary material are perfect substitutes in $R$ and therefore trade at the same price
$P^R$ per tonne. Treated waste trades at $P^T$ per tonne between the waste sector and the recycling
sector. Capital is rented at $P^K$ per unit per unit time by final-goods producers and recyclers
alike, so a single rental market allocates $K$ between $K^Y$ and $K^R$. The final good is bought at
its own price of one for consumption, for gross investment, and as an input to extraction, exploration
and waste treatment.

\smalltitle{Who hands over what.} The physical waste flow reaches the collection system through the
agent whose activity generates it, following the process-level accounting of
Section~\ref{sec:sp:setup:materialaccounting} line by line. It is useful to split the total into two
bases, because the instruments below treat them differently:
\begin{align}
W^{\mathcal D} &\equiv \Omega\mathcal D
= \Omega\left(\omega^YY+\omega^NC^N+\omega^DC^D+\omega^WC^W+\omega^CC\right),
&
W^{\circ} &\equiv \delta M^K+\mathcal N,
\label{eq:mkt:bases}
\end{align}
so that $W=W^{\mathcal D}+W^{\circ}$ by~\eqref{eq:sp:ledger-streams}. We call $W^{\mathcal D}$
\emph{use waste} --- material diffused when a final good is used up, handed over by the final-goods
producer ($\Omega^YY$), the miner ($\Omega^NC^N+\Omega^DC^D$), the waste firm itself ($\Omega^WC^W$)
and the household ($\Omega^CC$) --- and $W^{\circ}$ \emph{primary waste}: the material released by the
household's depreciating capital, $\delta M^K$, plus the mass the miner displaces from nature,
$\mathcal N=\Omega^{N,S}N+\Omega^{D,S}D$. Recycling hands over nothing of its own, since it draws no
final good and the depreciation of $K^R$ is already inside $\delta M^K$.

\subsection{Policy instruments}\label{sec:mkt:setup:instruments}

The government has seven instruments, in the three blocks anticipated above. All are stated as
time-varying rates and all revenue is returned to the household through a lump-sum transfer $\mathcal T$.

\smalltitle{(A) The emission block.} A tax $\tau^U$ per tonne of \emph{untreated} waste released to the
environment, paid by the waste sector; a tax $\tau^T$ per tonne of waste \emph{treated}, also paid by
the waste sector; and a subsidy $s^R$ per tonne of secondary material \emph{recovered}, paid to the
recycling sector. The three are one instrument in disguise. Since
$\Xi=\left(1-\varpi\right)W+d^W\left(\varpi W-R^R\right)$ by~\eqref{eq:sp:setup:Xi}, setting
\begin{align}
\tau^U &= \tau^\Xi,
&
\tau^T &= d^W\tau^\Xi,
&
s^R &= d^W\tau^\Xi
\label{eq:mkt:emission-split}
\end{align}
collects exactly $\tau^\Xi\Xi$: a uniform tax at rate $\tau^\Xi$ per tonne of waste reaching the
environment, with the liability split between the sector that controls the treated share and the
sector that controls the recovery rate. The split is not cosmetic. The waste firm chooses $\varpi$ but
cannot choose $a$, and the recycler chooses $a$ but does not release the residue; charging each of them
for the margin it actually controls is what makes a single emission price work through two sectors.

\smalltitle{(B) The materials block.} A tax $\tau^R$ per tonne of raw material delivered to final-goods
production, levied on the buyer, and a rebate $s^I$ per unit of \emph{gross} investment $G=I+\delta K$.
Together they form a deposit--refund scheme on materials: every tonne drawn into production pays
$\tau^R$ as though it were about to become waste, and the fraction of it that is instead locked up in
new durable capital is refunded. The rebate is the only instrument in the list that is levied on a
final-good flow rather than on a tonnage, and Section~\ref{sec:mkt:implementation} shows that its rate
has to be $\tau^R$ times the material content of investment goods, so it is a tonnage instrument after
all.

\smalltitle{(C) The waste-charge block.} A charge $\tau^{\mathcal D}$ per tonne of use waste and a
charge $\tau^{\circ}$ per tonne of primary waste, in both cases paid by the agent handing the waste
over, together with a subsidy $s^W$ per tonne collected, paid to the waste sector. The two charges are
gate fees at the point of collection; $s^W$ is what the public tender bids for. We do not impose
$\tau^{\mathcal D}=\tau^{\circ}$ in advance. Section~\ref{sec:mkt:implementation} shows that the two
rates must differ, and why.

Instruments the model does \emph{not} need are as instructive as the ones it does. There is no
consumption tax, no output tax, no extraction tax and no exploration tax, although the corresponding
wedges $1+\zeta\Omega^C$, $1-\zeta\Omega^Y$, $1+\zeta\Omega^N$ and $1+\zeta\Omega^D$ all appear in the
planner's conditions. The gate fee $\tau^{\mathcal D}$ delivers all four at once, because the thing
those four wedges have in common is a tonnage of use waste, and a charge per tonne is the natural base
for it. We return to this in Section~\ref{sec:mkt:implementation:reading}.

\subsection{Firms}\label{sec:mkt:setup:firms}

\smalltitle{Final goods.} The representative producer rents $K^Y$, buys $R$ at $P^R+\tau^R$ per tonne,
and hands over $\Omega^YY$ tonnes of use waste at the gate fee $\tau^{\mathcal D}$. It takes $P$ as
given. Its instantaneous problem is static,
\begin{align}
\max_{K^Y,R\ge0}\ \left(1-\tau^{\mathcal D}\Omega^Y\right)F\left(K^Y,R,P\right)-P^KK^Y-\left(P^R+\tau^R\right)R,
\label{eq:mkt:final-problem}
\end{align}
with first-order conditions
\begin{align}
F_K\left(1-\tau^{\mathcal D}\Omega^Y\right) &= P^K,
&
F_R\left(1-\tau^{\mathcal D}\Omega^Y\right) &= P^R+\tau^R .
\label{eq:mkt:final-foc}
\end{align}
Both are the familiar marginal-product conditions with output valued net of the gate fee its own waste
attracts. Note that the producer perceives no consequence of the material \emph{embodied} in the good
it sells: it pays for the tonnes it buys and for the tonnes it scraps, and the rest leaves with the
good. That is precisely the ledger externality of Section~\ref{sec:mkt:setup:externalities}, and
$\tau^R$ is what it costs to close it.

\smalltitle{Mining.} The mining firm owns the reserve, discounts at the household's rate $r$, and
solves an optimal-control problem in $\left(S,X\right)$. It pays gate fees at rate
$\tau^{\mathcal D}$ on the use waste its effort generates and at rate $\tau^{\circ}$ on the mass it
displaces from nature:
\begin{align}
\max_{N,D\ge0}\ \int_0^\infty
\Big[&P^RN-C^N\!\left(N,S\right)\left(1+\tau^{\mathcal D}\Omega^N\right)
-C^D\!\left(D,X\right)\left(1+\tau^{\mathcal D}\Omega^D\right)
\notag\\
&-\tau^{\circ}\left(\Omega^{N,S}N+\Omega^{D,S}D\right)\Big]
e^{-\int_0^{t}r\left(z\right)dz}\,dt,
\qquad
\dot S=D-N,
\quad
\dot X=D .
\label{eq:mkt:mining-problem}
\end{align}
Writing $P^S$ and $P^X$ for the firm's own shadow prices of $S$ and $X$, in final goods per tonne, the
conditions are
\begin{align}
P^R &= C^N_N\left(1+\tau^{\mathcal D}\Omega^N\right)+P^S+\tau^{\circ}\Omega^{N,S},
&
P^S+P^X &= C^D_D\left(1+\tau^{\mathcal D}\Omega^D\right)+\tau^{\circ}\Omega^{D,S},
\label{eq:mkt:mining-foc}
\end{align}
each with the corner branches of~\eqref{eq:sp:sum:N-corner} and~\eqref{eq:sp:sum:D-corner}, and
\begin{align}
\dot P^S &= rP^S-\left|C^N_S\right|\left(1+\tau^{\mathcal D}\Omega^N\right),
&
\dot P^X &= rP^X+C^D_X\left(1+\tau^{\mathcal D}\Omega^D\right),
\label{eq:mkt:mining-costates}
\end{align}
with the transversality conditions that select the forward solutions. These are
term-for-term the planner's~\eqref{eq:sp:sum:S} and~\eqref{eq:sp:sum:X} with
$\tau^{\mathcal D}$ in the place of $\zeta$: a private reserve owner already faces Hotelling's rule
with a stock effect, and the only thing policy has to supply is the price of the waste that mining
drags along with it.

\smalltitle{Recycling.} The recycler buys treated waste at $P^T$, rents capital at $P^K$, sells
secondary material at $P^R$, and collects the recovery subsidy $s^R$:
\begin{align}
\max_{T,K^R\ge0}\ \left(P^R+s^R\right)\mathcal R\!\left(T,K^R\right)-P^TT-P^KK^R .
\label{eq:mkt:recycler-problem}
\end{align}
With $\mathcal R_T=\alpha$ and $\mathcal R_{K^R}=a'$ from~\eqref{eq:sp:setup:recycling-derivatives},
and with $K^R$ bounded below exactly as in Section~\ref{sec:sp:opt:foc}, the conditions are
\begin{subequations}
\label{eq:mkt:recycler-foc}
\begin{align}
\alpha\left(P^R+s^R\right) &= P^T,
\label{eq:mkt:recycler-T}
\end{align}
and, for the capital margin,
\begin{align}
K^R &= 0,
&&\text{if}\quad a'\!\left(0^{+}\right)\left(P^R+s^R\right)\le P^K,
\label{eq:mkt:recycler-K-corner}
\\
a'\!\left(x\right)\left(P^R+s^R\right) &= P^K,
&&\text{if}\quad a'\!\left(0^{+}\right)\left(P^R+s^R\right)> P^K .
\label{eq:mkt:recycler-K-interior}
\end{align}
\end{subequations}
Because $\mathcal R$ is homogeneous of degree one, the two conditions exhaust the revenue,
$\left(P^R+s^R\right)\mathcal R=\left(P^R+s^R\right)\left(\alpha T+a'K^R\right)=P^TT+P^KK^R$, so the
recycling sector earns zero profit and pays no dividend. The sector's entire economic content is the
pair of prices at which it is willing to operate.

\smalltitle{Waste management.} The waste firm collects all $W$, which it takes as given, and chooses
only $\varpi$. It receives the collection subsidy $s^W$ per tonne and the price $P^T$ per tonne
treated, pays collection and treatment costs together with the gate fee on its own use waste, and pays
the two emission-block taxes:
\begin{align}
\Pi^W &= \Big[s^W+P^T\varpi
-\left(c^c+c^T\!\left(\varpi\right)\right)\left(1+\tau^{\mathcal D}\Omega^W\right)
-\tau^U\left(1-\varpi\right)-\tau^T\varpi\Big]W .
\label{eq:mkt:waste-profit}
\end{align}
Maximising over $\varpi\in\left[0,1\right]$ gives the three-branch condition
\begin{subequations}
\label{eq:mkt:waste-foc}
\begin{align}
\varpi &= 0
&&\text{if}\quad P^T+\tau^U-\tau^T\le0,
\\
P^T+\tau^U-\tau^T &= c^{T\prime}\!\left(\varpi\right)\left(1+\tau^{\mathcal D}\Omega^W\right)
&&\text{if}\quad 0<P^T+\tau^U-\tau^T<c^{T\prime}\!\left(1\right)\left(1+\tau^{\mathcal D}\Omega^W\right),
\\
\varpi &= 1
&&\text{if}\quad P^T+\tau^U-\tau^T\ge c^{T\prime}\!\left(1\right)\left(1+\tau^{\mathcal D}\Omega^W\right),
\end{align}
\end{subequations}
and the tender drives profit to zero, which determines the subsidy
\begin{align}
s^W &= c^W\left(1+\tau^{\mathcal D}\Omega^W\right)+\tau^U\left(1-\varpi\right)+\tau^T\varpi-P^T\varpi,
\qquad c^W\equiv c^c+c^T\!\left(\varpi\right),
\label{eq:mkt:waste-zero-profit}
\end{align}
evaluated at the $\varpi$ chosen in~\eqref{eq:mkt:waste-foc}. The firm treats waste for exactly two
reasons: treated waste can be sold, and treating it converts a tonne from the high tax rate $\tau^U$ to
the low rate $\tau^T$. The second motive survives when the first does not, which is what will keep
treatment alive in a phase where nothing is recycled. If $P^T$ rises far enough, $s^W$ turns negative
and the tender becomes a licence fee --- the right to collect waste is then worth paying for.

\subsection{The household}\label{sec:mkt:setup:household}

The household owns $K$ and, with it, the material embodied in the capital stock, $M^K$. It takes as
given the prices, the policy rates, the pollution stock $P$, and the economy-wide intensities
$\Omega^C$ and $\Phi^I$. It buys consumption at one plus its gate fee, buys gross investment at one
less the rebate, and pays the primary-waste charge on the material its capital releases as it
depreciates:
\begin{align}
\left(1+\tau^{\mathcal D}\Omega^C\right)C+\left(1-s^I\right)G+\tau^{\circ}\delta M^K
&= P^KK+\Pi+\mathcal T,
\label{eq:mkt:hh-budget}
\end{align}
where $\Pi$ collects the dividends of the firms it owns. Its two states evolve as
\begin{align}
\dot K &= G-\delta K,
&
\dot M^K &= \Phi^IG-\delta M^K .
\label{eq:mkt:hh-states}
\end{align}
Note the asymmetry that drives everything on this margin: the household chooses $G$ but not $\Phi^I$,
so a unit of gross investment brings with it a fixed tonnage of embodied material and, with it, a
future stream of primary-waste charges. That stream is a private liability, which is why $M^K$ is a
private state in the market economy as well as a social one.

Substituting~\eqref{eq:mkt:hh-budget} for $G$ into the current-value Hamiltonian with costates $\mu^K$
on $K$ and $\mu^M$ on $M^K$, and defining
\begin{align}
\lambda &\equiv \frac{\mu^K+\mu^M\Phi^I}{1-s^I},
&
p^M &\equiv -\frac{\mu^M}{\lambda},
&
q &\equiv \frac{\mu^K}{\lambda}\;=\;1-s^I+\Phi^Ip^M,
\label{eq:mkt:hh-costate-def}
\end{align}
so that $\lambda$ is the marginal utility of a final good spent, $p^M$ the private shadow cost of a
tonne locked into the capital stock and $q$ the price of installed capital, the conditions are
\begin{align}
\frac{u'\left(C\right)}{\lambda} &= 1+\tau^{\mathcal D}\Omega^C,
&
\dot p^M &= \left(r+\delta\right)p^M-\delta\tau^{\circ},
&
r &= \frac{P^K}{q}-\delta+\frac{\dot q}{q},
\label{eq:mkt:hh-foc}
\end{align}
with $r\equiv\rho-\dot\lambda/\lambda$ and the usual transversality conditions. The three have direct
readings. The first says the household equates marginal utility to the full price of a unit of
consumption, gate fee included. The second says the private cost of embodied material is the
discounted stream of the primary-waste charges it will attract as the capital runs off, at the
effective rate $r+\delta$ because the liability itself depreciates --- structurally identical to the
planner's~\eqref{eq:sp:costate:M}, with $\tau^{\circ}$ where the planner had $p^W$. The third is
ordinary capital-market arbitrage at the installed price $q$, and by~\eqref{eq:mkt:hh-costate-def} that
price departs from unity for two reasons the household can name: the rebate it received, and the
material liability it took on.

\subsection{Equilibrium}\label{sec:mkt:setup:equilibrium}

\begin{definition}[Competitive equilibrium]\label{def:mkt:equilibrium}
Given initial states $\left(K_0,S_0,X_0,P_0,M^K_0\right)$ and a policy
$\left\{\tau^U,\tau^T,s^R,\tau^R,s^I,\tau^{\mathcal D},\tau^{\circ}\right\}_{t\ge0}$, a competitive
equilibrium is a path of quantities $\left\{C,G,N,D,\varpi,K^R,W,\Omega\right\}$, prices
$\left\{P^R,P^T,P^K,\lambda,s^W\right\}$ and private shadow prices
$\left\{q,p^M,P^S,P^X\right\}$ such that: each agent solves its problem,
namely~\eqref{eq:mkt:final-foc}, \eqref{eq:mkt:mining-foc}--\eqref{eq:mkt:mining-costates},
\eqref{eq:mkt:recycler-foc}, \eqref{eq:mkt:waste-foc}--\eqref{eq:mkt:waste-zero-profit}
and~\eqref{eq:mkt:hh-budget}--\eqref{eq:mkt:hh-foc}; the four markets clear,
\begin{align}
R &= N+\mathcal R\!\left(T,K^R\right),
&
T &= \varpi W,
&
K^Y+K^R &= K,
\notag
\\
Y &= C+G+C^N+C^D+C^W,
&&&&
G &= I+\delta K;
\label{eq:mkt:clearing}
\end{align}
the materials accounting of Section~\ref{sec:sp:setup:aggregatematerial} holds,
\begin{align}
W &= \Omega\mathcal D+\delta M^K+\mathcal N,
&
R &= \Omega\left(\mathcal D+\phi^IG\right);
\label{eq:mkt:accounting}
\end{align}
the states obey~\eqref{eq:mkt:hh-states} together with $\dot S=D-N$, $\dot X=D$ and
$\dot P=\Xi-\theta\left(P\right)P$; and the government budget balances through $\mathcal T$.
\end{definition}

The count works out. Given the five states and the four private costates
$\left(q,p^M,P^S,P^X\right)$, the thirteen conditions consisting of the consumption
rule, the two final-goods conditions, the two mining conditions, the two recycling conditions, the
waste-treatment condition, the zero-profit condition, the household budget, the treated-waste market
clearing and the two accounting relations~\eqref{eq:mkt:accounting} determine the thirteen unknowns
$\left\{C,G,N,D,\varpi,K^R,W,\Omega,P^R,P^T,P^K,\lambda,s^W\right\}$, with $R$, $T$, $K^Y$ and $Y$
eliminated by definitions and the goods market clearing by Walras' law. The costates then propagate by
the third of~\eqref{eq:mkt:hh-foc}, the second of~\eqref{eq:mkt:hh-foc} and~\eqref{eq:mkt:mining-costates}.

Two differences from the planner's system in Section~\ref{sec:sp:opt:foc} are worth marking, because
they are the whole of what the government has to repair. First, the market economy carries \emph{four}
private costates against the planner's five: there is no $p^P$, because $P$ is nobody's state. Second,
where the planner carried $p^W$ and $\zeta$ as multipliers on the two accounting
relations~\eqref{eq:sp:ledger} and~\eqref{eq:sp:intensity}, the market carries neither. The
relations~\eqref{eq:mkt:accounting} still hold in equilibrium --- they are identities, not choices ---
but they hold without anyone being charged for them. Three shadow prices of the planned economy thus
have no private counterpart, and the policy of Section~\ref{sec:mkt:implementation} supplies exactly
three independent rates to stand in for them.
